\begin{figure}[htbp]
  \centering
  \resizebox{0.95\textwidth}{!}{%
  \begin{tikzpicture}[
    >=Latex,
    font=\small,
    node distance=20mm and 14mm,
    outerbox/.style={
      draw,
      rounded corners=3pt,
      minimum width=78mm,
      minimum height=50mm,
      fill=gray!6
    },
    innerbox/.style={
      draw,
      rounded corners=3pt,
      minimum width=60mm,
      align=left
    },
    caller/.style={
      draw,
      rounded corners=3pt,
      minimum width=30mm,
      align=center,
      fill=blue!6
    },
    arrow/.style={-{Latex[length=2.3mm]}, thick}
  ]
    \node[caller] (user) {Other code};

    \node[outerbox, right=24mm of user] (melody) {};
    \node[anchor=north] at (melody.north) {\textbf{Melody object}};

    \node[innerbox, fill=green!8, anchor=north] (public) at ([yshift=-6mm]melody.north) {
      \textbf{Public operations}\\
      \quad add\_event()\\
      \quad transpose()\\
      \quad to\_lilypond()
    };

    \node[innerbox, fill=yellow!10, anchor=south] (internal) at ([yshift=2mm]melody.south) {
      \textbf{Internal details}\\
      \quad events\\
      \quad cached length\\
      \quad harmony references
    };

    \draw[arrow] (user) -- node[above] {uses} (public.west);
  \end{tikzpicture}%
  }
  \caption{Illustration of encapsulation and abstraction. Encapsulation groups related data and behaviour inside one object, while abstraction lets other code use a small public interface without needing direct access to the internal representation.}
  \label{fig:oop_encapsulation_abstraction}
\end{figure}
